\documentclass[11pt,letterpaper]{article}
\usepackage[technical-reference]{cornell-notes}
\usepackage{needspace}

\title{Clause 18: Concepts Library - Cornell Notes}
\author{}
\date{}
\setDocTitle{Clause 18: Concepts Library}
\setDocSubtitle{C++ 2024 Cornell Notes}
\setDocOwner{C++ 2024 Cornell Notes}
\setDocAuthor{}
\setDocDate{}
\setCornellCollection{C++ 2024 Cornell Notes}
\setCornellUnitType{Clause}
\setCornellUnitNumber{18}
\setCornellUnitTitle{Concepts Library}
\setCornellCompanionLabel{C++ Technical Reference}
\hypersetup{pdftitle={Clause 18: Concepts Library - Cornell Notes},pdfsubject={C++ technical study notes},pdfauthor={}}

\begin{document}
\maketitle
\begin{CornellReferenceOrientationBox}
\textbf{Purpose.} Defines foundational language-support concepts for equality, ordering, object construction and movement, invocation, relations, and regularity. These concepts state semantic expectations used by generic library components.
\par\textbf{Role.} The vocabulary used to constrain and document generic algorithms and ranges.
\par\textbf{Principal dependencies.} Uses the library-wide rules of Clause 16 together with the applicable core-language concepts and supporting library clauses.
\par\textbf{Primary audience.} generic-library users and authors.
\par\textbf{Major subject groups.} General; Equality preservation; Header <concepts> synopsis; Language-related concepts; Comparison concepts; Object concepts; Callable concepts.
\end{CornellReferenceOrientationBox}
\section{Learning Objectives}
\begin{itemize}[label=\textcolor{CornellReferenceTeal}{$\blacktriangleright$}]
\item Define the governing ideas of General and apply its stated contract at a call site.
\item Identify the governing ideas of Equality preservation and apply its stated contract at a call site.
\item Distinguish the governing ideas of Header <concepts> synopsis and apply its stated contract at a call site.
\item Explain the governing ideas of Language-related concepts and apply its stated contract at a call site.
\item Trace the governing ideas of Comparison concepts and apply its stated contract at a call site.
\item Apply the governing ideas of Object concepts and apply its stated contract at a call site.
\item Analyze the governing ideas of Callable concepts and apply its stated contract at a call site.
\item Evaluate the governing ideas of General and apply its stated contract at a call site.
\item Diagnose the governing ideas of Concept same\_\allowbreak{}as and apply its stated contract at a call site.
\item Compare the governing ideas of Concept derived\_\allowbreak{}from and apply its stated contract at a call site.
\end{itemize}
\section{Normative Structure Map}
\small\begin{longtable}{@{}>{\RaggedRight\arraybackslash}p{0.13\textwidth}>{\RaggedRight\arraybackslash}p{0.27\textwidth}>{\RaggedRight\arraybackslash}p{0.19\textwidth}>{\RaggedRight\arraybackslash}p{0.31\textwidth}@{}}
\toprule\textbf{Subclause} & \textbf{Subject} & \textbf{Rule category} & \textbf{Key dependencies / entities}\\\midrule\endfirsthead
\toprule\textbf{Subclause} & \textbf{Subject} & \textbf{Rule category} & \textbf{Key dependencies / entities}\\\midrule\endhead
\bottomrule\endfoot
18.1 & General & library semantics & Uses the library-wide rules of Clause 16 together with the applicable core-language concepts and supporting library clauses.\\
18.2 & Equality preservation & library semantics & Uses the library-wide rules of Clause 16 together with the applicable core-language concepts and supporting library clauses.\\
18.3 & Header <concepts> synopsis & interface / declarations & Uses the library-wide rules of Clause 16 together with the applicable core-language concepts and supporting library clauses.\\
18.4 & Language-related concepts & library semantics & General, Concept same\_\allowbreak{}as, Concept derived\_\allowbreak{}from, Concept convertible\_\allowbreak{}to, and related subclauses\\
18.5 & Comparison concepts & library semantics & General, Boolean testability, Comparison common types, Concept equality\_\allowbreak{}comparable, and related subclauses\\
18.6 & Object concepts & library semantics & Uses the library-wide rules of Clause 16 together with the applicable core-language concepts and supporting library clauses.\\
18.7 & Callable concepts & library semantics & General, Concept invocable, Concept regular\_\allowbreak{}invocable, Concept predicate, and related subclauses\\
\end{longtable}\normalsize
\begin{CornellReferenceNormativeBox}A heading maps the clause; it does not replace the detailed conditions. For each operation, read requirements, effects, result, exceptions, complexity, remarks, and notes with their stated normative force.\end{CornellReferenceNormativeBox}
\subsection*{Rule Classification Guide}
\small\begin{longtable}{@{}>{\RaggedRight\arraybackslash}p{0.25\textwidth}>{\RaggedRight\arraybackslash}p{0.67\textwidth}@{}}\toprule\textbf{Label or category} & \textbf{How to read it}\\\midrule\endfirsthead\toprule\textbf{Label or category} & \textbf{How to read it (continued)}\\\midrule\endhead\bottomrule\endfoot
Constraints & Conditions controlling whether a declaration, overload, or specialization participates in the relevant context.\\
Mandates & Requirements checked for the described declaration or specialization; failure has the stated well-formedness consequence.\\
Preconditions & Obligations the caller establishes before evaluation; violating one forfeits the operation's contract.\\
Effects / Returns / Postconditions & Separate the state change, produced result, and state guaranteed after successful completion.\\
Throws / Error conditions & State how failure is reported and which exceptional paths are permitted or required.\\
Complexity & Bounds or exact counts for the operations named by the specification.\\
Remarks / Recommended practice & Remarks refine applicability or participation; recommended practice guides implementations without silently becoming a program requirement.\\
Exposition-only & A device for describing semantics, not a public name on which a program can depend.\\
\end{longtable}\normalsize
\section{Cornell Cue-and-Notes}
\Needspace{8\baselineskip}
\subsection*{18.1 General}
\begin{longtable}{@{}>{\RaggedRight\arraybackslash\columncolor{CornellReferenceCueGray}}p{0.28\textwidth}>{\RaggedRight\arraybackslash}p{0.64\textwidth}@{}}
\rowcolor{CornellReferenceNavy}\color{white}\textbf{Cue / recall prompt} & \color{white}\textbf{Notes}\\\endfirsthead
\rowcolor{CornellReferenceNavy}\color{white}\textbf{Cue / recall prompt} & \color{white}\textbf{Notes (continued)}\\\endhead
\arrayrulecolor{CornellReferenceLineGray}
\textbf{What contract governs General?}\\[-1mm]{\scriptsize\CornellClauseRef{18.1}} & Frames the terminology, applicability, and relationships needed to interpret General; detailed rules remain in the following subclauses.\\\hline
\end{longtable}
\begin{CornellReferenceSummaryBox}General supplies the library semantics portion of this clause. The key review move is to connect its local wording to the clause-wide model, then classify each condition by normative force before relying on it.\end{CornellReferenceSummaryBox}
\Needspace{8\baselineskip}
\subsection*{18.2 Equality preservation}
\begin{longtable}{@{}>{\RaggedRight\arraybackslash\columncolor{CornellReferenceCueGray}}p{0.28\textwidth}>{\RaggedRight\arraybackslash}p{0.64\textwidth}@{}}
\rowcolor{CornellReferenceNavy}\color{white}\textbf{Cue / recall prompt} & \color{white}\textbf{Notes}\\\endfirsthead
\rowcolor{CornellReferenceNavy}\color{white}\textbf{Cue / recall prompt} & \color{white}\textbf{Notes (continued)}\\\endhead
\arrayrulecolor{CornellReferenceLineGray}
\textbf{What contract governs Equality preservation?}\\[-1mm]{\scriptsize\CornellClauseRef{18.2}} & Specifies the facility family Equality preservation. Read its declarations together with mandates, preconditions, effects, results, exceptions, complexity, and lifetime rules.\\\hline
\end{longtable}
\begin{CornellReferenceSummaryBox}Equality preservation supplies the library semantics portion of this clause. The key review move is to connect its local wording to the clause-wide model, then classify each condition by normative force before relying on it.\end{CornellReferenceSummaryBox}
\Needspace{8\baselineskip}
\subsection*{18.3 Header <concepts> synopsis}
\begin{longtable}{@{}>{\RaggedRight\arraybackslash\columncolor{CornellReferenceCueGray}}p{0.28\textwidth}>{\RaggedRight\arraybackslash}p{0.64\textwidth}@{}}
\rowcolor{CornellReferenceNavy}\color{white}\textbf{Cue / recall prompt} & \color{white}\textbf{Notes}\\\endfirsthead
\rowcolor{CornellReferenceNavy}\color{white}\textbf{Cue / recall prompt} & \color{white}\textbf{Notes (continued)}\\\endhead
\arrayrulecolor{CornellReferenceLineGray}
\textbf{What contract governs Header <concepts> synopsis?}\\[-1mm]{\scriptsize\CornellClauseRef{18.3}} & Catalogs the declarations made available by Header <concepts> synopsis. The synopsis is an interface map; the detailed specifications supply constraints and behavior.\\\hline
\end{longtable}
\begin{CornellReferenceSummaryBox}Header <concepts> synopsis supplies the interface / declarations portion of this clause. The key review move is to connect its local wording to the clause-wide model, then classify each condition by normative force before relying on it.\end{CornellReferenceSummaryBox}
\Needspace{8\baselineskip}
\subsection*{18.4 Language-related concepts}
\begin{longtable}{@{}>{\RaggedRight\arraybackslash\columncolor{CornellReferenceCueGray}}p{0.28\textwidth}>{\RaggedRight\arraybackslash}p{0.64\textwidth}@{}}
\rowcolor{CornellReferenceNavy}\color{white}\textbf{Cue / recall prompt} & \color{white}\textbf{Notes}\\\endfirsthead
\rowcolor{CornellReferenceNavy}\color{white}\textbf{Cue / recall prompt} & \color{white}\textbf{Notes (continued)}\\\endhead
\arrayrulecolor{CornellReferenceLineGray}
\textbf{What contract governs General?}\\[-1mm]{\scriptsize\CornellClauseRef{18.4.1}} & Frames the terminology, applicability, and relationships needed to interpret General; detailed rules remain in the following subclauses.\\\hline
\textbf{What contract governs Concept same\_\allowbreak{}as?}\\[-1mm]{\scriptsize\CornellClauseRef{18.4.2}} & Defines or applies the generic requirement Concept same\_\allowbreak{}as; syntactic satisfaction must be paired with every stated semantic and equality-preservation expectation.\\\hline
\textbf{What contract governs Concept derived\_\allowbreak{}from?}\\[-1mm]{\scriptsize\CornellClauseRef{18.4.3}} & Defines or applies the generic requirement Concept derived\_\allowbreak{}from; syntactic satisfaction must be paired with every stated semantic and equality-preservation expectation.\\\hline
\textbf{What contract governs Concept convertible\_\allowbreak{}to?}\\[-1mm]{\scriptsize\CornellClauseRef{18.4.4}} & Defines or applies the generic requirement Concept convertible\_\allowbreak{}to; syntactic satisfaction must be paired with every stated semantic and equality-preservation expectation.\\\hline
\textbf{What contract governs Concept common\_\allowbreak{}reference\_\allowbreak{}with?}\\[-1mm]{\scriptsize\CornellClauseRef{18.4.5}} & Defines or applies the generic requirement Concept common\_\allowbreak{}reference\_\allowbreak{}with; syntactic satisfaction must be paired with every stated semantic and equality-preservation expectation.\\\hline
\textbf{What contract governs Concept common\_\allowbreak{}with?}\\[-1mm]{\scriptsize\CornellClauseRef{18.4.6}} & Defines or applies the generic requirement Concept common\_\allowbreak{}with; syntactic satisfaction must be paired with every stated semantic and equality-preservation expectation.\\\hline
\textbf{What contract governs Arithmetic concepts?}\\[-1mm]{\scriptsize\CornellClauseRef{18.4.7}} & Defines or applies the generic requirement Arithmetic concepts; syntactic satisfaction must be paired with every stated semantic and equality-preservation expectation.\\\hline
\textbf{What contract governs Concept assignable\_\allowbreak{}from?}\\[-1mm]{\scriptsize\CornellClauseRef{18.4.8}} & Defines or applies the generic requirement Concept assignable\_\allowbreak{}from; syntactic satisfaction must be paired with every stated semantic and equality-preservation expectation.\\\hline
\textbf{What contract governs Concept swappable?}\\[-1mm]{\scriptsize\CornellClauseRef{18.4.9}} & Governs state transfer or exchange for Concept swappable; check self-operation, validity after movement, exception behavior, and invalidation.\\\hline
\textbf{What contract governs Concept destructible?}\\[-1mm]{\scriptsize\CornellClauseRef{18.4.10}} & Defines or applies the generic requirement Concept destructible; syntactic satisfaction must be paired with every stated semantic and equality-preservation expectation.\\\hline
\textbf{What contract governs Concept constructible\_\allowbreak{}from?}\\[-1mm]{\scriptsize\CornellClauseRef{18.4.11}} & Defines or applies the generic requirement Concept constructible\_\allowbreak{}from; syntactic satisfaction must be paired with every stated semantic and equality-preservation expectation.\\\hline
\textbf{What contract governs Concept default\_\allowbreak{}initializable?}\\[-1mm]{\scriptsize\CornellClauseRef{18.4.12}} & Determines the initialization form used by Concept default\_\allowbreak{}initializable, the candidates or conversions considered, object-lifetime start, narrowing rules, and failure consequences.\\\hline
\textbf{What contract governs Concept move\_\allowbreak{}constructible?}\\[-1mm]{\scriptsize\CornellClauseRef{18.4.13}} & Defines or applies the generic requirement Concept move\_\allowbreak{}constructible; syntactic satisfaction must be paired with every stated semantic and equality-preservation expectation.\\\hline
\textbf{What contract governs Concept copy\_\allowbreak{}constructible?}\\[-1mm]{\scriptsize\CornellClauseRef{18.4.14}} & Defines or applies the generic requirement Concept copy\_\allowbreak{}constructible; syntactic satisfaction must be paired with every stated semantic and equality-preservation expectation.\\\hline
\end{longtable}
\begin{CornellReferenceSummaryBox}Language-related concepts supplies the library semantics portion of this clause. The key review move is to connect its local wording to the clause-wide model, then classify each condition by normative force before relying on it.\end{CornellReferenceSummaryBox}
\Needspace{8\baselineskip}
\subsection*{18.5 Comparison concepts}
\begin{longtable}{@{}>{\RaggedRight\arraybackslash\columncolor{CornellReferenceCueGray}}p{0.28\textwidth}>{\RaggedRight\arraybackslash}p{0.64\textwidth}@{}}
\rowcolor{CornellReferenceNavy}\color{white}\textbf{Cue / recall prompt} & \color{white}\textbf{Notes}\\\endfirsthead
\rowcolor{CornellReferenceNavy}\color{white}\textbf{Cue / recall prompt} & \color{white}\textbf{Notes (continued)}\\\endhead
\arrayrulecolor{CornellReferenceLineGray}
\textbf{What contract governs General?}\\[-1mm]{\scriptsize\CornellClauseRef{18.5.1}} & Frames the terminology, applicability, and relationships needed to interpret General; detailed rules remain in the following subclauses.\\\hline
\textbf{What contract governs Boolean testability?}\\[-1mm]{\scriptsize\CornellClauseRef{18.5.2}} & Specifies the facility family Boolean testability. Read its declarations together with mandates, preconditions, effects, results, exceptions, complexity, and lifetime rules.\\\hline
\textbf{What contract governs Comparison common types?}\\[-1mm]{\scriptsize\CornellClauseRef{18.5.3}} & Specifies the facility family Comparison common types. Read its declarations together with mandates, preconditions, effects, results, exceptions, complexity, and lifetime rules.\\\hline
\textbf{What contract governs Concept equality\_\allowbreak{}comparable?}\\[-1mm]{\scriptsize\CornellClauseRef{18.5.4}} & Defines or applies the generic requirement Concept equality\_\allowbreak{}comparable; syntactic satisfaction must be paired with every stated semantic and equality-preservation expectation.\\\hline
\textbf{What contract governs Concept totally\_\allowbreak{}ordered?}\\[-1mm]{\scriptsize\CornellClauseRef{18.5.5}} & Defines or applies the generic requirement Concept totally\_\allowbreak{}ordered; syntactic satisfaction must be paired with every stated semantic and equality-preservation expectation.\\\hline
\end{longtable}
\begin{CornellReferenceSummaryBox}Comparison concepts supplies the library semantics portion of this clause. The key review move is to connect its local wording to the clause-wide model, then classify each condition by normative force before relying on it.\end{CornellReferenceSummaryBox}
\Needspace{8\baselineskip}
\subsection*{18.6 Object concepts}
\begin{longtable}{@{}>{\RaggedRight\arraybackslash\columncolor{CornellReferenceCueGray}}p{0.28\textwidth}>{\RaggedRight\arraybackslash}p{0.64\textwidth}@{}}
\rowcolor{CornellReferenceNavy}\color{white}\textbf{Cue / recall prompt} & \color{white}\textbf{Notes}\\\endfirsthead
\rowcolor{CornellReferenceNavy}\color{white}\textbf{Cue / recall prompt} & \color{white}\textbf{Notes (continued)}\\\endhead
\arrayrulecolor{CornellReferenceLineGray}
\textbf{What contract governs Object concepts?}\\[-1mm]{\scriptsize\CornellClauseRef{18.6}} & Defines or applies the generic requirement Object concepts; syntactic satisfaction must be paired with every stated semantic and equality-preservation expectation.\\\hline
\end{longtable}
\begin{CornellReferenceSummaryBox}Object concepts supplies the library semantics portion of this clause. The key review move is to connect its local wording to the clause-wide model, then classify each condition by normative force before relying on it.\end{CornellReferenceSummaryBox}
\Needspace{8\baselineskip}
\subsection*{18.7 Callable concepts}
\begin{longtable}{@{}>{\RaggedRight\arraybackslash\columncolor{CornellReferenceCueGray}}p{0.28\textwidth}>{\RaggedRight\arraybackslash}p{0.64\textwidth}@{}}
\rowcolor{CornellReferenceNavy}\color{white}\textbf{Cue / recall prompt} & \color{white}\textbf{Notes}\\\endfirsthead
\rowcolor{CornellReferenceNavy}\color{white}\textbf{Cue / recall prompt} & \color{white}\textbf{Notes (continued)}\\\endhead
\arrayrulecolor{CornellReferenceLineGray}
\textbf{What contract governs General?}\\[-1mm]{\scriptsize\CornellClauseRef{18.7.1}} & Frames the terminology, applicability, and relationships needed to interpret General; detailed rules remain in the following subclauses.\\\hline
\textbf{What contract governs Concept invocable?}\\[-1mm]{\scriptsize\CornellClauseRef{18.7.2}} & Defines or applies the generic requirement Concept invocable; syntactic satisfaction must be paired with every stated semantic and equality-preservation expectation.\\\hline
\textbf{What contract governs Concept regular\_\allowbreak{}invocable?}\\[-1mm]{\scriptsize\CornellClauseRef{18.7.3}} & Defines or applies the generic requirement Concept regular\_\allowbreak{}invocable; syntactic satisfaction must be paired with every stated semantic and equality-preservation expectation.\\\hline
\textbf{What contract governs Concept predicate?}\\[-1mm]{\scriptsize\CornellClauseRef{18.7.4}} & Defines or applies the generic requirement Concept predicate; syntactic satisfaction must be paired with every stated semantic and equality-preservation expectation.\\\hline
\textbf{What contract governs Concept relation?}\\[-1mm]{\scriptsize\CornellClauseRef{18.7.5}} & Defines or applies the generic requirement Concept relation; syntactic satisfaction must be paired with every stated semantic and equality-preservation expectation.\\\hline
\textbf{What contract governs Concept equivalence\_\allowbreak{}relation?}\\[-1mm]{\scriptsize\CornellClauseRef{18.7.6}} & Defines or applies the generic requirement Concept equivalence\_\allowbreak{}relation; syntactic satisfaction must be paired with every stated semantic and equality-preservation expectation.\\\hline
\textbf{What contract governs Concept strict\_\allowbreak{}weak\_\allowbreak{}order?}\\[-1mm]{\scriptsize\CornellClauseRef{18.7.7}} & Defines or applies the generic requirement Concept strict\_\allowbreak{}weak\_\allowbreak{}order; syntactic satisfaction must be paired with every stated semantic and equality-preservation expectation.\\\hline
\end{longtable}
\begin{CornellReferenceSummaryBox}Callable concepts supplies the library semantics portion of this clause. The key review move is to connect its local wording to the clause-wide model, then classify each condition by normative force before relying on it.\end{CornellReferenceSummaryBox}
\section{Language-Lawyer and Library-Usage Pitfalls}
\begin{CornellReferencePitfallBox}\begin{itemize}[label=\textcolor{CornellReferenceAmber}{$\blacktriangleright$}]
\item Reading a synopsis as the complete behavioral contract.
\item Ignoring mandates, preconditions, exception behavior, or complexity requirements.
\item Using an exposition-only name as though it were public API.
\item Assuming invalidation, ownership, or thread-safety properties that are not stated.
\item Treating successful compilation as proof that semantic requirements and preconditions are met.
\end{itemize}\end{CornellReferencePitfallBox}
\section{Programmer and Implementation Checklist}
\begin{CornellReferenceChecklistBox}\begin{itemize}[label=$\square$]
\item Identify the declaring header or module exposure for each facility.
\item Check template arguments and mandates before evaluating an operation.
\item Establish every precondition at the call site.
\item Record effects, returned value, postconditions, and exceptions separately.
\item Audit iterator, reference, pointer, and object lifetime after mutation.
\item Verify stated complexity and synchronization assumptions.
\item For \CornellClauseRef{18.1}, verify general against the precise rule category and all cited dependencies.
\item For \CornellClauseRef{18.2}, verify equality preservation against the precise rule category and all cited dependencies.
\item For \CornellClauseRef{18.3}, verify header <concepts> synopsis against the precise rule category and all cited dependencies.
\item For \CornellClauseRef{18.4}, verify language-related concepts against the precise rule category and all cited dependencies.
\end{itemize}\end{CornellReferenceChecklistBox}
\section{Clause Synthesis}
\begin{CornellReferenceOrientationBox}Defines foundational language-support concepts for equality, ordering, object construction and movement, invocation, relations, and regularity. These concepts state semantic expectations used by generic library components. The vocabulary used to constrain and document generic algorithms and ranges. A sound review distinguishes syntax and participation from runtime preconditions, keeps notes and examples non-mandatory, and follows cross-clause dependencies before making a portability claim.\end{CornellReferenceOrientationBox}
\section{Self-Test Questions}
\begin{enumerate}
\item What role does General play in the clause's overall model?
\item Which normative distinctions must be kept separate when applying Equality preservation?
\item How would you audit a program or implementation that depends on Header <concepts> synopsis?
\item Compare Language-related concepts with Comparison concepts; what different question does each answer?
\item What role does Comparison concepts play in the clause's overall model?
\item Which normative distinctions must be kept separate when applying Object concepts?
\item How would you audit a program or implementation that depends on Callable concepts?
\item Compare General with Concept same\_\allowbreak{}as; what different question does each answer?
\item Scenario: a program relies on an unstated implementation behavior while using Concept same\_\allowbreak{}as. What must be checked before calling the program portable?
\item Scenario: a program relies on an unstated implementation behavior while using Concept derived\_\allowbreak{}from. What must be checked before calling the program portable?
\end{enumerate}
\clearpage\section{Answer Key}
\begin{enumerate}
\item General is one part of the clause's library semantics model. Its role is understood by connecting the local rule in 18.7.1 to the clause orientation and its dependent concepts.
\item Keep formation and well-formedness separate from runtime preconditions and effects. Also distinguish mandatory wording from notes, implementation choice from undefined behavior, and interface declarations from detailed contracts; see 18.2.
\item Audit Header <concepts> synopsis by locating the controlling wording in 18.3, listing inputs and type constraints, proving any preconditions, recording effects and failure behavior, and checking lifetime, invalidation, complexity, and synchronization where applicable.
\item Language-related concepts and Comparison concepts occupy different steps or facility groups. Analyze each under its own subclause, then follow the explicit dependency between them rather than transferring one set of guarantees to the other; begin at 18.4.
\item Comparison concepts is one part of the clause's library semantics model. Its role is understood by connecting the local rule in 18.5 to the clause orientation and its dependent concepts.
\item Keep formation and well-formedness separate from runtime preconditions and effects. Also distinguish mandatory wording from notes, implementation choice from undefined behavior, and interface declarations from detailed contracts; see 18.6.
\item Audit Callable concepts by locating the controlling wording in 18.7, listing inputs and type constraints, proving any preconditions, recording effects and failure behavior, and checking lifetime, invalidation, complexity, and synchronization where applicable.
\item General and Concept same\_\allowbreak{}as occupy different steps or facility groups. Analyze each under its own subclause, then follow the explicit dependency between them rather than transferring one set of guarantees to the other; begin at 18.7.1.
\item First identify the exact requirement in 18.4.2, then classify the relied-on behavior as required, unspecified, implementation-defined, conditionally supported, or undefined. Portability requires satisfying the program obligations and avoiding dependence on a choice not guaranteed in the target environments.
\item First identify the exact requirement in 18.4.3, then classify the relied-on behavior as required, unspecified, implementation-defined, conditionally supported, or undefined. Portability requires satisfying the program obligations and avoiding dependence on a choice not guaranteed in the target environments.
\end{enumerate}
\section{Key-Term Recap}
\begin{longtable}{@{}>{\RaggedRight\arraybackslash}p{0.28\textwidth}>{\RaggedRight\arraybackslash}p{0.64\textwidth}@{}}\toprule\textbf{Term} & \textbf{Working meaning}\\\midrule\endfirsthead\toprule\textbf{Term} & \textbf{Working meaning (continued)}\\\midrule\endhead\bottomrule\endfoot
same\_\allowbreak{}as & A symmetric concept requiring two types to be the same.\\
convertible\_\allowbreak{}to & A concept combining implicit and explicit conversion expectations.\\
regular\_\allowbreak{}invocable & An invocable whose call is required to preserve equality and not modify function-object or argument state through the invocation.\\
predicate & A Boolean-testable regular invocation.\\
strict\_\allowbreak{}weak\_\allowbreak{}order & A relation satisfying irreflexive and transitive ordering properties.\\
regular & A semiregular type that is equality comparable.\\
\end{longtable}
\CornellTakeaway{Concept satisfaction includes semantic expectations, not merely whether a requires-expression compiles.}
\end{document}
